\chapter*{Vorwort: Das schwierigste Gebiet}
\addcontentsline{toc}{chapter}{Vorwort}
\markboth{Vorwort}{Vorwort}

Die Statistik ist der einzige Teil der Mathematik, den fast jeder anwendet und fast niemand versteht. Das ist keine Polemik, sondern eine empirisch gut belegte Aussage --- und der Grund, warum dieser Band anders gebaut ist als seine Vorgänger.

\begin{merksatz}
\boxtitel{Zwei Befunde, die am Anfang stehen sollten}
\begin{center}
\begin{minipage}{0.9\textwidth}
\textbf{Erstens.} Befragungen unter Wissenschaftlerinnen und Wissenschaftlern, die regelmäßig Statistik anwenden, zeigen seit Jahrzehnten dasselbe Bild: Ein erheblicher Teil kann den $p$-Wert nicht korrekt definieren. In einer vielzitierten Untersuchung an Medizinerinnen und Medizinern wählte die Mehrheit mindestens eine falsche von sechs vorgelegten Deutungen als richtig.

\textbf{Zweitens.} Dasselbe gilt für das Konfidenzintervall. Die verbreitete Deutung --- \enquote{der wahre Wert liegt mit $95\,\%$ Wahrscheinlichkeit in diesem Intervall} --- ist falsch, und zwar nicht in einem spitzfindigen Sinn.
\end{minipage}
\end{center}

\begin{center}
\begin{minipage}{0.88\textwidth}
\kursiv{Beide Begriffe werden in diesem Band vollständig hergeleitet, und beide bekommen einen eigenen Warnkasten, der die falschen Deutungen einzeln durchgeht. Wenn Sie nach Kapitel~6 und Kapitel~7 sagen können, was ein $p$-Wert und ein Konfidenzintervall \kursiv{nicht} sind, hat der Band seinen Zweck erfüllt.}
\end{minipage}
\end{center}
\end{merksatz}

\section*{Warum das so ist}

Es liegt nicht an mangelnder Sorgfalt. Es liegt daran, dass die Statistik zwei Dinge vermischt, die in der übrigen Mathematik getrennt bleiben.

\begin{merksatz}
\boxtitel{Die zwei Hälften}
\begin{center}
\renewcommand{\arraystretch}{1.55}
\begin{tabular}{@{}P{2.6cm}P{6.0cm}@{}}
\toprule
\textbf{Die Rechnung} & Ist gewöhnliche Mathematik. Sie lässt sich beweisen, und dieser Band beweist sie.\\
\textbf{Das Urteil} & Ob ein Verfahren im vorliegenden Fall \kursiv{anwendbar} ist, folgt aus keiner Rechnung. Es ist eine Entscheidung über die Wirklichkeit.\\
\bottomrule
\end{tabular}
\end{center}

\begin{center}
\begin{minipage}{0.88\textwidth}
\kursiv{In Band 16 konnte man sich auf die Definitionen zurückziehen: Ein Graph \kursiv{ist} ein Paar $(V,E)$, basta. Hier steht vor jeder Rechnung die Frage, ob die Voraussetzungen zutreffen --- und diese Frage beantwortet die Statistik nicht, sondern das Fach, aus dem die Daten kommen.}
\end{minipage}
\end{center}
\end{merksatz}

\begin{warnung}
\boxtitelr{Die Folge: ein Band mit doppelter Buchführung}
Zu jedem Verfahren gehören in diesem Band drei Dinge:

\begin{center}
\renewcommand{\arraystretch}{1.5}
\begin{tabular}{@{}P{2.4cm}P{6.2cm}@{}}
\toprule
\textbf{Die Herleitung} & Woher die Formel kommt, vollständig nachgerechnet.\\
\textbf{Die Voraussetzungen} & Was gelten muss, damit sie gilt --- ausdrücklich aufgezählt.\\
\textbf{Der Bruchfall} & Was geschieht, wenn eine Voraussetzung verletzt ist. Meist ein Rechenbeispiel.\\
\bottomrule
\end{tabular}
\end{center}

\begin{center}
\begin{minipage}{0.88\textwidth}
\kursiv{Der dritte Punkt fehlt in den meisten Lehrbüchern, und er ist der wichtigste. Ein Verfahren, dessen Bruchverhalten man nicht kennt, wendet man blind an.}
\end{minipage}
\end{center}
\end{warnung}

\section*{Was Sie mitbringen müssen}

\begin{merksatz}
\boxtitel{Die Voraussetzungen}
\begin{center}
\renewcommand{\arraystretch}{1.5}
\begin{tabular}{@{}P{3.4cm}P{5.2cm}@{}}
\toprule
\textbf{Schulrechnen} & Bruchrechnung, Wurzeln, Potenzen, Umformen von Gleichungen.\\
\textbf{Summenzeichen} & $\sum$ als Abkürzung für \enquote{addiere alles auf}. Wird in Kapitel~2 noch einmal ausführlich eingeführt.\\
\textbf{Wahrschein\-lichkeit} & Nur die Grundbegriffe, und sie werden in Kapitel~4 wiederholt. Band 4 der Reihe enthält alles ausführlich.\\
\textbf{Integrale} & Erst ab Kapitel~4, und dort nur als Flächenberechnung gedeutet. Wer sie nicht kennt, kommt über die Tabellen durch.\\
\bottomrule
\end{tabular}
\end{center}

\begin{center}
\begin{minipage}{0.88\textwidth}
\kursiv{Dieser Band setzt Band 4 \enquote{Die Vermessung des Zufalls} nicht voraus, verweist aber häufig dorthin. Wer die Wahrscheinlichkeitstheorie dahinter sehen will, findet sie dort; wer nur rechnen will, kommt ohne sie aus.}
\end{minipage}
\end{center}
\end{merksatz}

\section*{Die MOBIL-Studie}

Dieser Band hat eine Besonderheit, die die anderen nicht haben: ein einziges Beispiel, das durch alle fünfzehn Kapitel läuft.

\begin{merksatz}
\boxtitel{Die Studie, die uns begleitet}
\begin{center}
\begin{minipage}{0.9\textwidth}
An einer orthopädischen Klinik wird geprüft, ob ein neues Schmerzschema nach Hüftgelenksersatz die Patientinnen und Patienten schneller wieder auf die Beine bringt als das bisherige.

$120$ Personen werden zufällig auf zwei Gruppen verteilt:
\end{minipage}
\end{center}

\begin{center}
\renewcommand{\arraystretch}{1.5}
\begin{tabular}{@{}P{2.0cm}P{2.0cm}P{4.0cm}@{}}
\toprule
\textbf{Gruppe A} & $60$ Personen & bisheriges Schema\\
\textbf{Gruppe B} & $60$ Personen & neues Schema\\
\bottomrule
\end{tabular}
\end{center}

\schritt{Was erhoben wird}
\begin{center}
{\small
\renewcommand{\arraystretch}{1.45}
\begin{tabular}{@{}P{3.4cm}P{2.4cm}P{2.6cm}@{}}
\toprule
\textbf{Merkmal} & \textbf{Kurzname} & \textbf{Werte}\\
\midrule
Schmerz nach $24$ h & \code{NRS24} & $0$ bis $10$\\
Zeit bis zur Mobilisation & \code{MOB} & Stunden\\
Zusätzliche Schmerzmittel & \code{BEDARF} & ja / nein\\
Zufriedenheit & \code{ZUFR} & $1$ bis $5$\\
Alter & \code{ALTER} & Jahre\\
Geschlecht & \code{GESCHL} & w / m\\
Operationsdauer & \code{OPDAUER} & Minuten\\
\bottomrule
\end{tabular}}
\end{center}

\begin{center}
\begin{minipage}{0.88\textwidth}
\kursiv{Die Studie ist erfunden, die Zahlen sind es auch. Sie sind aber so gewählt, dass sie sich von Hand rechnen lassen und dass die Ergebnisse realistisch ausfallen. Jedes Verfahren dieses Bandes wird an ihr vorgeführt.}
\end{minipage}
\end{center}
\end{merksatz}

\begin{merksatz}
\boxtitel{Warum ein durchgehendes Beispiel}
\begin{center}
\renewcommand{\arraystretch}{1.5}
\begin{tabular}{@{}P{3.0cm}P{5.6cm}@{}}
\toprule
\textbf{Vergleich\-barkeit} & Verschiedene Verfahren auf dieselben Daten angewandt geben verschiedene Antworten. Das sieht man nur, wenn die Daten dieselben bleiben.\\
\textbf{Der ganze Weg} & Von der Häufigkeitstabelle über den Test bis zur Metaanalyse --- an einem einzigen Vorhaben.\\
\textbf{Die Fallen bleiben sichtbar} & Wer in Kapitel~2 den Mittelwert des NRS gebildet hat, erfährt in Kapitel~9, warum das fragwürdig war.\\
\bottomrule
\end{tabular}
\end{center}
\end{merksatz}

\section*{Drei Warnungen}

\begin{warnung}
\boxtitelr{Erstens: Signifikanz ist nicht Bedeutsamkeit}
Das Wort \enquote{signifikant} bedeutet in der Statistik etwas anderes als in der Umgangssprache. Es heißt: \kursiv{Ein Ergebnis dieser Art wäre unter der Nullhypothese selten.} Es heißt nicht: wichtig, groß, klinisch bedeutsam oder gar bewiesen.

Bei genügend großer Stichprobe wird jeder noch so winzige Unterschied signifikant. Kapitel~7 rechnet das vor.
\end{warnung}

\begin{warnung}
\boxtitelr{Zweitens: Kein Test beweist die Nullhypothese}
Ein nicht signifikantes Ergebnis heißt \kursiv{nicht}, dass kein Unterschied besteht. Es heißt nur, dass die Daten nicht ausreichen, um einen nachzuweisen.

\enquote{Wir konnten keinen Unterschied zeigen} und \enquote{Es gibt keinen Unterschied} sind verschiedene Aussagen. Die Verwechslung ist der häufigste Fehler in medizinischen Veröffentlichungen.
\end{warnung}

\begin{warnung}
\boxtitelr{Drittens: Die Rechnung ist der kleinste Teil}
Jedes Verfahren dieses Bandes lässt sich in einer Zeile aufrufen. Die Rechnung dauert Millisekunden.

Die Arbeit steckt vorher --- in der Frage, welches Verfahren passt --- und nachher --- in der Frage, was das Ergebnis bedeutet. Ein Lehrbuch, das nur die Rechnung zeigt, bereitet auf den unwichtigen Teil vor.

\begin{center}
\begin{minipage}{0.88\textwidth}
\kursiv{Aus demselben Grund enthält dieser Band fast keinen Programmcode. Wo er nützlich ist, steht ein kurzer Technikkasten mit zwei bis drei Zeilen R --- zur Kontrolle der Handrechnung, nicht als Ersatz für sie.}
\end{minipage}
\end{center}
\end{warnung}

\section*{Der Aufbau}

\begin{merksatz}
\boxtitel{Sechs Teile, fünfzehn Kapitel}
\begin{center}
{\small
\renewcommand{\arraystretch}{1.45}
\begin{tabular}{@{}P{1.0cm}P{3.2cm}P{4.4cm}@{}}
\toprule
\textbf{Teil} & \textbf{Thema} & \textbf{Die Leitfrage}\\
\midrule
I & Beschreiben (1--3) & Was steht in den Daten?\\
II & Schätzen (4--6) & Was sagen sie über die Grundgesamtheit?\\
III & Testen (7--9) & Wann ist ein Unterschied mehr als Zufall?\\
IV & Modelle (10--11) & Wie hängen mehrere Größen zusammen?\\
V & Medizin (12--14) & Was ist beim Menschen anders?\\
VI & Stand (15) & Was ist heute strittig?\\
\bottomrule
\end{tabular}}
\end{center}
\end{merksatz}

\vspace{10pt}
\begin{flushright}
\itshape\color{grau}
Die Formel ist das Leichte.\\
Die Frage, ob sie hier gilt,\\
ist das Eigentliche.
\end{flushright}
